\bstctlcite{IEEEtran:BSTcontrol}
\maketitle

%% ====================================================================
\begin{abstract}
\ifqceshort
We present a hardware-aware hybrid pipeline that feeds learned Hetionet
knowledge-graph (KG) RotatE embeddings into a 16-qubit Pauli quantum-kernel
classifier within a stacking ensemble for \textit{Compound-treats-Disease}
(CtD) link prediction (47{,}031 entities, 2.25 M edges).
On the reproducible protocol (cached 256D RotatE, MoA features, Nystr\"{o}m
$m{=}200$, HistGBDT-led stacking), we report \textbf{0.7805} test PR-AUC
(seed 42) and \textbf{0.7398 $\pm$ 0.038} over five seeds; QSVC remains near
chance, so we do not claim quantum lift on CtD.
A legacy \texttt{--fast\_mode} run showed Pauli-QSVC errors decorrelate from
classical tree models (+5.8 pp ensemble under that protocol).
Clinical validation reveals \textit{score-validity inversion}; a ten-feature
MoA module and Tier-3 RNA-seq repurposing validation (CREEDS fusion) address
downstream ranking honesty.
\else
Repurposing approved drugs for new indications bypasses the most expensive
stages of \textit{de novo} discovery and can deliver therapies in 3--6 years
instead of the customary 10--15 year, \$2--3 billion pipeline.
We present a hardware-aware hybrid pipeline that feeds learned knowledge-graph
(KG) vector embeddings into quantum kernel classifiers for biomedical link
prediction, targeting \textit{Compound-treats-Disease} (CtD) relationships on
the Hetionet KG (47{,}031 entities, 2.25 M edges, 24 relation types).
Full-graph RotatE training maps each entity into a 128-dimensional space;
pairs are classified by a stacking ensemble whose base learners include tuned
Random Forest, Extra Trees models, and a quantum support vector classifier
(QSVC) built from a 16-qubit Pauli fidelity kernel.
A key finding is that the Pauli feature map---whose multi-axis entanglement
produces a Hilbert-space geometry qualitatively different from tree-based
classifiers---contributes non-redundant signal that raises ensemble
PR-AUC by $+5.8$ pp despite lowering standalone QSVC accuracy by $-8.7$ pp.
Clinical validation against ClinicalTrials.gov exposes a
\textit{score-validity inversion}: the highest-scoring predictions are not
the best-supported clinically.
We introduce a ten-feature mechanism-of-action (MoA) module---binding-target
overlap, pathway co-involvement, pharmacologic class, and treatment
analogy---to re-align model scores with clinical validity.
On the current reproducible protocol (cached 256D RotatE, MoA features, full
classical suite, Nystr\"{o}m $m{=}200$), the stacking ensemble reaches
\textbf{0.7805} test PR-AUC (seed 42) and \textbf{0.7398 $\pm$ 0.038} over five
seeds; we do not claim quantum advantage on CtD under this protocol.
\fi
\end{abstract}

\medskip\noindent\textbf{Keywords:} quantum machine learning; quantum kernel methods; knowledge graph embeddings;
link prediction; drug repurposing; quantum support vector machine; stacking
ensemble; Hetionet.

%% ====================================================================
\section{Introduction}
%% ====================================================================

Drug repurposing offers a faster, lower-cost path to clinical deployment than
\textit{de novo} discovery because candidate molecules already carry established
safety and pharmacokinetic data.
Biomedical knowledge graphs (KGs) such as Hetionet \cite{himmelstein2017}
aggregate heterogeneous evidence from genomics, pharmacology, pathways, and
clinical observation, encoding 2.25 M relationships across 47,031 entities.
The \textit{Compound-treats-Disease} (CtD) relation (755 known pairs) defines
a well-posed link-prediction task.

KG vector embedding methods (TransE, RotatE, ComplEx) project entities into a
continuous latent space in which relational structure is approximately preserved.
In RotatE, $\mathbf{h}\circ\mathbf{r}\approx\mathbf{t}$, so training on all
24 Hetionet relation types enriches compound and disease embeddings with
multi-hop context inaccessible to purely molecular descriptors.

Quantum machine learning offers kernel methods computing similarity in
exponentially large Hilbert spaces.
How prior quantum drug-discovery work relates to molecular inputs, and how our
use of learned KG embeddings differs, is spelled out in
Section~\ref{sec:critical_distinction}.
To our knowledge, no prior work feeds learned KG embeddings into quantum feature
spaces for link prediction; this paper bridges the two research tracks.

\subsection{Contributions}

\begin{enumerate}[leftmargin=*,label=\arabic*.]
\item \textbf{Novel intersection.} Hardware-aware integration of RotatE KG
  embeddings with quantum kernel classifiers (QSVC, fidelity kernel) for
  biomedical link prediction.
\item \textbf{Counterintuitive quantum-classical synergy (legacy protocol).}
  Pauli feature map degrades standalone QSVC ($-8.7$ pp) yet raised ensemble
  PR-AUC ($+5.8$ pp) under historical \texttt{--fast\_mode}; current multiseed
  runs show no quantum lift when HistGBDT dominates.
\item \textbf{MoA feature module.} Ten pharmacological-plausibility features
  from Hetionet structure address the score-validity inversion inherent in
  embedding-based scoring.
\item \textbf{Clinical validation.} Two of six novel predictions (Losartan
  $\to$ atherosclerosis, Mitomycin $\to$ liver cancer) are supported by 7+
  clinical trials each; Ezetimibe $\to$ gout is a biologically plausible novel
  hypothesis.
\item \textbf{Reproducible pipeline} with configurable embeddings, feature maps,
  dimensionality reduction, stacked ensembles, and IBM Quantum hardware
  backends.
\end{enumerate}

\ifqceshort
%% ====================================================================
\section{Background}
%% ====================================================================
Prior quantum drug-discovery work uses molecular fingerprints; classical KG
repurposing uses embeddings alone \cite{mayers2023}. We route learned RotatE
vectors from full Hetionet into a Pauli fidelity-kernel QSVC within a stacking
ensemble \cite{havlicek2019,schuld2019}.
\else
%% ====================================================================
\section{Background}
%% ====================================================================

\subsection{Hetionet and KG-Based Drug Repurposing}

Hetionet v1.0 \cite{himmelstein2017} has 47,031 nodes and 2,250,198 edges
across 24 relation types.
Key relations: CtD (755), Compound-binds-Gene (CbG, 11,571), Disease-associates-Gene
(DaG, 12,623), Gene-participates-Pathway (GpPW, 84,372), Compound-resembles-Compound
(CrC, 6,486), Disease-resembles-Disease (DrD, 543).
Mayers \emph{et al.}\ \cite{mayers2023} achieved MRR 0.9792 via ensemble on
biomedical KGs; recent work integrates LLM retrieval \cite{llmrag2025} and
multi-database KGs \cite{quantum2026}, all classical.

\subsection{Quantum Kernel Methods}

Quantum kernels compute pairwise similarity in a quantum feature space:
$k(\mathbf{x},\mathbf{y})=|\langle 0|U^\dagger(\mathbf{x})U(\mathbf{y})|0\rangle|^2$,
where $U(\mathbf{x})$ is a parameterized encoding circuit.
QSVC uses this kernel as a drop-in replacement for classical kernels
\cite{havlicek2019,schuld2019}.
Kernel computation scales as $O(n^2)$ in training size \cite{schuld2021}.

\subsection{Stacking Ensembles and Model Diversity}

Stacking trains a meta-learner on base-model predictions, exploiting diverse
errors rather than individually high performance \cite{wolpert1992}.
Hybrid quantum-classical ensembles add a quantum base learner whose errors are
weakly correlated with classical ones
\cite{havlicek2019,schuld2019,liu2021,abbas2021,incudini2023}.

\subsection{Prior quantum drug discovery vs.\ KG embeddings}
\label{sec:critical_distinction}

\textbf{The critical distinction:} To the best of the authors' knowledge---and
we acknowledge we cannot have surveyed every concurrent group, preprint, or
adjacent line of work; the field moves quickly, and as an emerging research
team we state this modestly and welcome corrections---all prior quantum drug
discovery work applies quantum kernels or circuits to molecular-level
features---fingerprints, QSAR descriptors, binding affinities.
KG-based drug repurposing, meanwhile, has been pursued with classical methods
\cite{mayers2023,llmrag2025,quantum2026}.
Our contribution is to route \emph{learned} KG embedding vectors (RotatE over the
full Hetionet graph) into a quantum kernel classifier for link prediction, rather
than molecular fingerprints alone.
\fi

%% ====================================================================
\section{Dataset}
%% ====================================================================

\ifqceshort
CtD uses 755 positives with 1:1 hard negatives (80/20 split); classical
features are 1177D (1187D with MoA); quantum path uses 16 qubits after PCA.
\else
\begin{table}[t]
\caption{Hetionet CtD Dataset Statistics}
\label{tab:dataset}
\centering\small
\begin{tabular}{@{}ll@{}}
\toprule
\textbf{Statistic} & \textbf{Value} \\
\midrule
Total entities / edges & 47,031 / 2,250,198 \\
Target relation (CtD) & 755 positive edges \\
Train positives / negatives & 604 / 604 (1:1 hard neg.) \\
Test positives / negatives & 151 / 151 \\
Feature dim.\ (classical) & 1177 (base); 1187 (+MoA) \\
Feature dim.\ (quantum) & 16 qubits (1177D$\to$24D$\to$16Q) \\
\bottomrule
\end{tabular}
\end{table}

CtD positives are split 80/20; hard negatives use degree-weighted head/tail
corruption, rejecting known positives.
The low tail uniqueness (75 of 137 diseases in training) reflects Hetionet's
concentration on well-studied diseases.
\fi

\ifqceshort\else
%% ====================================================================
\section{Methods Overview}
%% ====================================================================

This work presents a hybrid quantum-classical machine learning pipeline for
compound-treats-disease link prediction using biomedical knowledge graph
embeddings derived from Hetionet. Node representations are first generated
using RotatE full-graph embeddings and augmented with topology-derived
similarity features including cosine similarity, Euclidean distance,
degree-based measures, and shortest-path structure.

Feature dimensionality is reduced to align with near-term quantum hardware
constraints before encoding into quantum feature maps including
\textit{PauliFeatureMap} and \textit{ZZFeatureMap}. These embeddings are
evaluated using fidelity-based quantum kernel estimation and classified with
quantum support vector classifiers.

Quantum models are compared against classical baselines including logistic
regression, random forest, and extra trees classifiers. Predictions are
combined through a stacking ensemble architecture to evaluate whether quantum
kernels contribute complementary geometric structure improving link prediction
robustness.

To address structural artifacts common in knowledge graph link prediction, a
mechanism-of-action correction module is introduced to distinguish
pharmacologically meaningful predictions from proximity-driven graph
correlations.
\fi

%% ====================================================================
\section{Hybrid Pipeline}
%% ====================================================================

\subsection{Pipeline Overview}

\begin{enumerate}[leftmargin=*,label=\arabic*.]
\item Load Hetionet; extract CtD edges as positives.
\item Train full-graph RotatE embeddings (all 24 relation types, PyKEEN \cite{pykeen2021}).
\item Construct train/test pairs with hard negative sampling
  (1:1 ratio; Section~\ref{sec:moa}).%
  \footnote{%
  \emph{Hard negatives} are constructed by \emph{corrupting} each positive
  compound--disease edge: with fixed probability we replace either the head or
  the tail by a \emph{degree-weighted} alternative entity of the same kind
  (compound vs.\ disease), rejecting pairs that coincide with a known positive
  CtD edge.
  This yields negatives that are graph-theoretically closer to true positives
  than uniform random non-edges, making the classification task harder and more
  realistic than i.i.d.\ random sampling.%
  }
\item Build pair feature vectors from embeddings and graph topology.
\item \textbf{Classical path:} Train RF, ET, LR with optional GridSearchCV.
\item \textbf{Quantum path:} PCA $\to$ Pauli/ZZ feature map $\to$ fidelity kernel $\to$ QSVC.
\item \textbf{Stacking ensemble:} Logistic regression meta-learner on OOF base probabilities.
\item Evaluate on held-out test set; report PR-AUC and ROC-AUC.
\end{enumerate}

\begin{figure}[htbp]
\centering
\includegraphics[width=\columnwidth]{../figures/fig1_pipeline.pdf}
\caption{Hybrid quantum-classical pipeline for CtD link prediction on Hetionet.
Full-graph RotatE embeddings feed a stacking ensemble combining tuned RF/ET/LR
(classical path) and a 16-qubit Pauli fidelity-kernel QSVC (quantum path).}
\label{fig:pipeline}
\end{figure}

\subsection{Full-Graph KG Embeddings}

RotatE \cite{rotate2019} models each relation as an element-wise complex
rotation: $\mathbf{h}\circ\mathbf{r}\approx\mathbf{t}$,
where $\mathbf{h},\mathbf{t}\in\mathbb{C}^d$ are the \emph{head} and
\emph{tail} entity embeddings (the two endpoints of a graph edge, e.g.\
a compound and a disease), $\mathbf{r}\in\mathbb{C}^d$ is the
\emph{relation} embedding encoding the edge type (e.g.\ \emph{treats}),
$\circ$ denotes element-wise multiplication, and $|\mathbf{r}_i|=1$ so
each component acts as a unit-modulus rotation in the complex plane.
Intuitively, a valid triple $(h, r, t)$---such as
\textit{Losartan}~$\xrightarrow{\mathrm{treats}}$~\textit{Atherosclerosis}---should
satisfy the constraint after training; invalid triples will not.
Training on all 2.25\,M edges (primary: 128D, 200 epochs; extended: 256D,
250 epochs) enriches compound/disease embeddings with gene-binding, pathway,
side-effect, and pharmacologic-class context unavailable to molecular descriptors.

\subsection{Pair Feature Construction}

For each $(c,d)$ pair, \texttt{EnhancedFeatureBuilder} produces a 1177-dim
vector $\mathbf{x}_{cd}$:
\textbf{(i)} Embedding block ($9d+3$ dims): concatenation, element-wise
$|\mathbf{h}-\mathbf{t}|$, Hadamard $\mathbf{h}\odot\mathbf{t}$, three global
similarity scalars, and polynomial features.
\textbf{(ii)} Graph topology (14 dims): degree, common neighbors, Jaccard,
Adamic--Adar, PageRank, betweenness, clustering---all computed on
\emph{training} edges only.
\textbf{(iii)} Domain features (8 dims): entity-type indicators and metaedge
counts.

\ifqceshort\else
\subsection{Classical Base Models}

Three base learners operate on $\mathbf{x}_{cd}$:
\textbf{RF} (500 estimators, Gini, GridSearchCV-tuned),
\textbf{ET} (500 estimators, fully randomised splits, same grid),
\textbf{LR} (L2, $C\in[10^{-3},10^2]$).
All produce calibrated probability outputs $p_m(y=1|\mathbf{x}_{cd})$.
\fi

\subsection{Quantum Kernel Classifier}

\textbf{Dimensionality reduction.} 1177D $\xrightarrow{\text{PCA}}$ 24D
$\xrightarrow{\text{linear}}$ 16 qubits (primary); direct $\to$ 12 qubits
(256D extended). Feature dimensionality was selected to remain compatible
with execution on near-term superconducting quantum hardware including IBM
Heron-class processors, ensuring feasibility beyond simulator-based
evaluation. Dimensionality reduction was performed not only for computational
efficiency but also to align embedding structure with qubit-limited quantum
kernel evaluation regimes typical of current NISQ devices.

\textbf{Feature maps.} Two Qiskit \cite{qiskit2023} maps are evaluated:
\emph{ZZ} (second-order Pauli-Z, reps=2) and
\emph{Pauli} (mixed X,Y,Z,ZZ, full entanglement, reps=2).

\textbf{Fidelity kernel.}
$k_Q(\mathbf{x},\mathbf{y})=|\langle 0^{\otimes n}|U^\dagger(\mathbf{x})U(\mathbf{y})|0^{\otimes n}\rangle|^2$,
computed via Qiskit Aer CPU statevector simulation \cite{qiskit_aer2023}.
The historical 128D headline run used no Nystr\"{o}m approximation
\cite{williams2001nystrom,qiskit_ml2023} (\texttt{nystrom\_m=None}); the full
$1{,}208\!\times\!1{,}208$ kernel matrix required \textbf{2,619 s}.
Reproducible multiseed benchmarks use Nystr\"{o}m landmarks
(\texttt{--qsvc\_nystrom\_m 200}).
\textbf{QSVC} uses $C\!=\!0.1$ (primary) or $C\!=\!0.676$ (Optuna-tuned);
$C$ is the standard SVM soft-margin parameter \cite{qiskit_ml2023}.
Rather than serving as a standalone classifier replacement, the quantum support
vector classifier contributes complementary geometric structure within the
ensemble architecture, improving prediction diversity and robustness relative
to purely classical feature spaces.

\subsection{Stacking Ensemble}

\ifqceshort
Level-0 learners (RF, ET, LR, QSVC) supply OOF probabilities to a logistic
regression meta-learner; decorrelated QSVC errors can lift ensemble PR-AUC
when classical bases dominate.
\else
Level-0 base learners (RF, ET, LR, QSVC) generate out-of-fold (OOF) probabilities
via $K=5$ stratified cross-validation. QSVC probabilities are obtained via
Platt scaling \cite{platt1999} (\texttt{probability=True} in \texttt{sklearn.svm.SVC},
inherited by QSVC \cite{qiskit_ml2023}):
\begin{equation}
\mathbf{z}_i=\bigl[p_\mathrm{RF}^\mathrm{OOF},p_\mathrm{ET}^\mathrm{OOF},
p_\mathrm{LR}^\mathrm{OOF},p_\mathrm{QSVC}^\mathrm{OOF}\bigr](\mathbf{x}_i).
\label{eq:stack}
\end{equation}
A level-1 logistic regression meta-learner maps $\mathbf{z}_i$ to the final
probability $\hat{p}(\mathbf{x}_i)=\sigma(\mathbf{w}^\top\mathbf{z}_i+b)$.
The variance of the weighted meta-prediction satisfies
\begin{equation}
\mathrm{Var}[\hat{p}]=\sum_{m,m'}w_m w_{m'}\,\mathrm{Cov}(\epsilon_m,\epsilon_{m'}),
\label{eq:ensemble_var}
\end{equation}
so weakly correlated base-model errors directly reduce ensemble variance.
Stacking allows the model to capture complementary inductive biases across
classical tree-based learners and quantum kernel similarity structure,
improving predictive stability across heterogeneous feature geometries.
\fi

\ifqceshort\else
\subsubsection*{Hyperparameter optimisation}
Optuna \cite{optuna2019} performs a two-stage TPE search: a quantum-structure
study caches each kernel matrix per $(\mathrm{map},\mathrm{reps},n_Q,\mathrm{PCA})$
tuple, then a classical/meta-learner study sweeps RF/ET/LR and stacking
parameters on the frozen kernel.
\fi

%% ====================================================================
\section{Mechanism-of-Action Feature Module}
\label{sec:moa}
%% ====================================================================

\ifqceshort
Ten MoA features from Hetionet structure (shared targets, pathways, analogies)
address score-validity inversion; activated with \texttt{--use\_moa\_features}.
\else
Graph-embedding scores measure structural proximity, which correlates with
treatment plausibility on average but fails in dense hub regions.
This correction step reduces false positives produced by hub-node proximity
effects that commonly inflate link prediction confidence in dense biomedical
interaction graphs.
The Abacavir $\to$ ocular cancer case (Section~\ref{sec:clinical}) illustrates
the failure: HIV comorbidity pulls the two entities close in embedding space
despite zero pharmacological rationale.

We compute ten features directly from Hetionet multi-relational structure for
every $(c,d)$ pair.
They are organised into three biological channels: \emph{direct target evidence}
(features 1--5), \emph{drug-class analogy} (features 6--8), and
\emph{disease-class analogy} (features 9--10).
All use only training-set CtD edges to avoid leakage.

The MoA module adds 10 features ($1177\to1187$ dims) and is activated with
\texttt{--use\_moa\_features}.
Its features include \texttt{shared\_targets}, \texttt{target\_overlap},
\texttt{shared\_pathway}, pharmacologic class, chemical similarity, and
similar-disease treatment analogies. On the training set,
\texttt{shared\_targets} shows an order-of-magnitude separation:
mean $\approx0.31$ for positives vs.\ $\approx0.04$ for hard negatives.

A natural question is whether the KG embedding itself could be retrained to
incorporate mechanistic plausibility directly --- e.g.\ by adding MoA-derived
supervision (a weighted CbG\,$\cap$\,DaG triple-loss head, or PCiC class
regularisation) to PyKEEN training, producing embeddings that natively
separate structural proximity from mechanistic plausibility and removing the
need for a post-hoc feature module. We treat the present MoA module as a
deliberate post-hoc correction that demonstrates the existence and structure
of the score-validity inversion problem.

All workflows were designed for compatibility with hardware execution
pipelines and reproducible comparison across simulator and quantum processor
environments.
Taken together, the pipeline represents a hardware-aware hybrid
quantum-classical inference workflow rather than a simulator-only evaluation
of quantum kernel classifiers.
\fi

%% ====================================================================
\section{Experimental Setup}
\label{sec:experimental}
%% ====================================================================

\textbf{Software.} Python 3.9+; PyKEEN 1.10+ (KG embeddings); Qiskit 1.x and
Qiskit-Aer (quantum circuits); scikit-learn 1.4+ (classical models, stacking);
PyTorch 2.x (PyKEEN backend).

\textbf{Hardware.} Qiskit Aer CPU statevector simulator (primary); IBM Quantum
Heron (hardware validation). Seed 42 for all splits and training.

\textbf{Reproducible Tier-1 benchmark (reported headline):}
{\small\begin{verbatim}
./scripts/run_256d_moa_multiseed.sh
\end{verbatim}}
Uses cached full-graph RotatE 256D embeddings (symlinked as 128D paths),
MoA features, full classical tuning (no \texttt{--fast\_mode}),
\texttt{--qsvc\_nystrom\_m 200}, and stacking that selects the best
PR-AUC classical base learner (HistGBDT) for the meta-learner
(\texttt{\_select\_best\_ensemble\_classical}).
\textbf{Caveat:} \texttt{--fast\_mode} trains only RF+ET and caps embedding
epochs at 50; it must not be used for HistGBDT/MoA benchmarking.
Historical single-run PR-AUC 0.7987 used \texttt{--fast\_mode} with
RF-led classical selection and is not reproduced on the current protocol
(see \texttt{results/headline\_repro\_fresh/INVESTIGATION.md}).
Results: \texttt{results/rerun\_256d\_moa/}; summary in
\texttt{results/multiseed/TABLE3.md}.

%% ====================================================================
\section{Results}
%% ====================================================================

We report PR-AUC as the primary metric (rare-positive class on the
underlying graph; invariant to negative-pool size, unlike ROC-AUC) and
ROC-AUC and accuracy as secondary metrics. Optuna optimises PR-AUC during
hyperparameter search. A random classifier on the 1:1 balanced test set
gives PR-AUC $\approx 0.50$; all reported models exceed this substantially.

\subsection{Primary Results}

\begin{table}[t]
\caption{Reproducible Tier-1 CtD benchmark (cached 256D RotatE + MoA, five seeds).
Nystr\"{o}m $m{=}200$; stacking uses best PR-AUC classical (HistGBDT).
Source: \texttt{results/rerun\_256d\_moa/summary.json}.}
\label{tab:primary}
\centering\small
\setlength{\tabcolsep}{3pt}
\begin{tabular}{@{}lc@{}}
\toprule
\textbf{Model} & \textbf{Test PR-AUC} \\
\midrule
\textbf{Ensemble-QC-stacking} & \textbf{0.7398 $\pm$ 0.038} \\
HistGBDT & 0.7393 $\pm$ 0.037 \\
\midrule
\multicolumn{2}{l}{\footnotesize Best single run (seed 42): ensemble \textbf{0.7805}; HistGBDT 0.7784.} \\
\bottomrule
\end{tabular}
\end{table}

Table~\ref{tab:primary} reports the reproducible headline; the full pipeline is shown in Figure~\ref{fig:pipeline}.
On seed 42, stacking adds +0.0021 PR-AUC over HistGBDT; QSVC alone remains near chance ($\approx$0.50) on this protocol, so we do not claim quantum lift on CtD.

\ifqceshort\else
\begin{table}[t]
\caption{Historical single-run Pauli inversion (128D, full kernel, \texttt{--fast\_mode}; not current headline protocol).}
\label{tab:legacy_primary}
\centering\small
\setlength{\tabcolsep}{3pt}
\begin{tabular}{@{}lccc@{}}
\toprule
\textbf{Model} & \textbf{PR-AUC} & \textbf{ROC-AUC} & \textbf{Acc.} \\
\midrule
Ensemble-QC (Pauli) & 0.7987 & 0.7456 & 0.576 \\
RandomForest-Opt & 0.7838 & 0.7319 & 0.583 \\
QSVC (Pauli) & 0.6343 & 0.6313 & 0.586 \\
\bottomrule
\end{tabular}
\end{table}

\subsection{Feature Map Analysis: The Pauli Inversion Effect}
\label{sec:results_feature_map}

\begin{table}[t]
\caption{Feature map ablation (RotatE-128D constant).}
\label{tab:ablation}
\centering\small
\begin{tabular}{@{}lcc@{}}
\toprule
\textbf{Feature Map} & \textbf{QSVC PR-AUC} & \textbf{Ensemble PR-AUC} \\
\midrule
ZZ (reps=2, 908~s kernel) & \textbf{0.7216} & 0.7408 \\
\textbf{Pauli (reps=2, 2,619~s)} & 0.6343 & \textbf{0.7987} \\
\bottomrule
\end{tabular}
\end{table}

\begin{figure}[htbp]
\centering
\includegraphics[width=\columnwidth]{../figures/fig2_pauli_zz.pdf}
\caption{Pauli vs.\ ZZ feature map: switching to Pauli lowers standalone
QSVC by 8.7 pp while raising ensemble PR-AUC by 5.8 pp.
Dashed line: best classical baseline (RF, 0.7838).}
\label{fig:pauli_zz}
\end{figure}

Figure~\ref{fig:pauli_zz} and Table~\ref{tab:ablation} show the central
quantum finding of this paper: the four-bar comparison fixes everything
(RotatE-128D, $C{=}0.1$, hard negatives, identical split) except the feature
map. Standalone, ZZ-QSVC scores 0.7216 and Pauli-QSVC scores 0.6343
($-8.7$~pp). \emph{In the ensemble, the order inverts:} the ZZ ensemble
reaches only 0.7408 (below RF at 0.7838), while the Pauli ensemble reaches
\textbf{0.7987} ($+5.8$~pp over ZZ ensemble; $+1.5$~pp over RF). The naive
``pick the higher single-model bar'' heuristic would have selected ZZ and
missed the headline result.

The mechanism is prediction \emph{decorrelation}. RF and ET split the
1177-dim feature vector along axis-aligned thresholds; the ZZ kernel's
restricted Pauli-$Z$ basis preserves much of that coordinate-aligned
structure, so ZZ-QSVC errors overlap with tree errors and the meta-learner
gains little. The Pauli map's mixed X, Y, Z and ZZ interactions across all
qubit pairs produce a 16-qubit Hilbert-space geometry qualitatively unlike
any axis-aligned partition, so its errors are decorrelated from RF/ET. The
meta-learner exploits this diversity: Pauli-QSVC is individually weaker but
\emph{more informative} as a base learner. The relevant ensemble objective
is therefore classical--quantum prediction decorrelation, not standalone
quantum accuracy. We caveat this with the usual single-seed, single-relation
qualification; multi-seed confirmation is a v2 deliverable.

\subsection{Hardware Validation: IBM Quantum Heron}

Hardware runs on a 16-qubit Pauli configuration yield QSVC PR-AUC
${\approx}0.634$, consistent with the simulator result (0.6343), confirming
that the fidelity kernel signal is observable within the gate-error budget on current NISQ
hardware.
\fi

%% ====================================================================
\section{Clinical Validation}
\label{sec:clinical}
%% ====================================================================

\ifqceshort
Top novel predictions were checked against ClinicalTrials.gov.
Losartan $\to$ atherosclerosis (score 0.528) and Mitomycin $\to$ liver cancer
(0.525) each have 7+ supporting trials; Abacavir $\to$ ocular cancer scores
highest (0.793) with zero trials---illustrating \textit{score-validity inversion}
that the MoA module targets.
\else
\begin{table}[t]
\caption{ClinicalTrials.gov validation of top novel predictions.}
\label{tab:clinical}
\centering\small
\setlength{\tabcolsep}{3pt}
\begin{tabularx}{\columnwidth}{@{}c >{\raggedright\arraybackslash}X c
  >{\raggedright\arraybackslash}p{0.16\columnwidth}
  >{\raggedright\arraybackslash}X@{}}
\toprule
\textbf{Rank} & \textbf{Prediction} & \textbf{Score} &
  \textbf{Trials} & \textbf{Verdict} \\
\midrule
1 & Abacavir $\to$ Ocular Cancer & 0.793 & 0 & Graph artifact \\
2 & Ezetimibe $\to$ Gout & 0.693 & 0 direct; 4 anti-inflam. & Plausible novel \\
3 & Ramipril $\to$ Stomach Ca. & 0.597 & 0 & No support \\
4 & Losartan $\to$ Atherosclerosis & 0.528 & 7+ (Ph.~4) & \textbf{Validated} \\
5 & Mitomycin $\to$ Liver Ca. & 0.525 & 7 (TACE) & \textbf{Validated} \\
6 & Salmeterol $\to$ Liver Ca. & 0.520 & 0 & No support \\
\bottomrule
\end{tabularx}
\end{table}

The highest-scoring prediction (Abacavir $\to$ ocular cancer, 0.793) has zero
clinical trial support while the two most validated predictions score only
0.52--0.53---a \textit{score-validity inversion}.
Root cause: RotatE embeddings absorb graph-structural proximity from HIV
comorbidity pathways, not mechanistic plausibility.
Abacavir's \texttt{shared\_targets}~=~0 and \texttt{target\_overlap}~=~0
would flag this case under the MoA module.

\textbf{Novel hypothesis.}
Ezetimibe $\to$ gout: NPC1L1 expression in proximal tubules relevant to urate
reabsorption, and emerging evidence linking cholesterol to urate transport,
provide biological plausibility despite zero direct trial support.

\begin{figure}[htbp]
\centering
\includegraphics[width=\columnwidth]{../figures/fig3_clinical.pdf}
\caption{Score-validity inversion: high model score need not imply clinical
support; lower-scoring Losartan and Mitomycin links have trial evidence.}
\label{fig:clinical}
\end{figure}
\fi

%% ====================================================================
\section{Discussion}
%% ====================================================================

\ifqceshort
Under a legacy \texttt{--fast\_mode} protocol, Pauli-QSVC errors were
decorrelated from RF/ET, raising ensemble PR-AUC despite weaker standalone
QSVC; the reproducible multiseed headline (Table~\ref{tab:primary}) shows no
quantum lift when HistGBDT dominates.
We use 16 qubits as a practical NISQ operating point on IBM Heron-class
hardware: enough expressivity for KG-derived features while limiting circuit
depth and $O(n^2)$ kernel cost on CtD's 755-edge scale.
\else
\subsection{Why the Ensemble Beats Classical-Only}

The Pauli quantum kernel provides +5.8 pp ensemble gain over ZZ despite
performing $-8.7$ pp worse in isolation
(Tables~\ref{tab:primary} and~\ref{tab:ablation}; Figure~\ref{fig:pauli_zz}).
RF/ET learn axis-aligned splits on the 1177-dim feature vector; the Pauli
kernel encodes similarity via 16-qubit Hilbert-space inner products.
The meta-learner weights QSVC as a ``second opinion from a different geometry.''
This decorrelation benefit is independent of absolute QSVC accuracy, suggesting
quantum kernel selection for ensembles should maximize
\emph{classical-quantum error independence}, not standalone quantum performance.

\subsection{Embedding Quality}

Embedding quality remains the dominant classical driver: moving from RotatE
128D to 256D raises the RF baseline by +7.3 pp and the ensemble ceiling by
+5.9 pp. This supports the NIER hypothesis that quantum kernels are most useful
when fed relation-rich KG embeddings rather than raw molecular descriptors.

\subsection{Quantum Kernel Scaling}

The fidelity kernel requires $O(n^2)$ evaluations: ${\approx}$1.46 M for the
primary 16-qubit configuration (2,619 s).
For larger relations (CbG: 11,571 edges $\to$ ${\sim}$18,500 pairs),
Nystr\"{o}m approximation ($m\approx100$ landmarks) is required.
The 755-edge CtD relation sits at the practical boundary for exact kernel
computation on current simulators.

We selected a 16-qubit feature encoding dimension to balance expressivity with
near-term hardware feasibility. This dimensionality allows fidelity-based
kernel evaluation while remaining compatible with current superconducting
quantum processors such as IBM Heron-class systems. Increasing the number of
qubits would increase representational capacity of the quantum feature space
and could improve kernel separability for higher-order relational structure,
but would also increase circuit depth, noise sensitivity, and execution cost.
The 16-qubit configuration therefore represents a practical operating point
between model expressiveness and NISQ-era hardware constraints.

\subsection{Score-Validity Inversion and MoA}

Pure embedding-based scoring rewards structural proximity
(Table~\ref{tab:clinical}); the MoA module
reintroduces the mechanistic axis via discrete set-operation counts
(\texttt{shared\_targets}, \texttt{target\_overlap}, \texttt{shared\_pathway\_genes})
that are orthogonal to RotatE geometry and interpretable by clinicians.
\fi

%% ====================================================================
\section{Conclusion}
%% ====================================================================

We have presented a hardware-aware hybrid quantum-classical pipeline bridging
KG embeddings with quantum kernel classifiers for biomedical link prediction.
On Hetionet CtD, our reproducible stacking ensemble achieves
\textbf{0.7805} test PR-AUC (256D RotatE + MoA, seed 42) and
\textbf{0.7398 $\pm$ 0.038} over five seeds (Table~\ref{tab:primary}).
The Pauli inversion effect\ifqceshort\else\ (Table~\ref{tab:legacy_primary})\fi\ reframes quantum
kernel selection: maximize classical-quantum prediction \emph{decorrelation},
not standalone accuracy, though current multiseed runs show no quantum lift
when HistGBDT dominates.
Clinical validation identifies Losartan $\to$ atherosclerosis and Mitomycin
$\to$ liver cancer as strongly supported predictions and reveals a
score-validity inversion motivating the MoA feature module.
RNA-seq repurposing validation (harmonized TCGA$\to$GSE225846; CREEDS fusion)
is documented in \texttt{docs/repurposing/RNASEQ\_REPURPOSING\_METHODS.md}.

\section*{Acknowledgment}
The authors thank Abdelrahman Elsayed, the IBM HBCU Quantum Center, and the IBM
Quantum Advocate Mentorship Program for mentorship and research support.

%% ====================================================================
\bibliographystyle{IEEEtran}
\bibliography{references_qce26}

